% chapters/ch04_intertemporal_consumption_and_asset_pricing.tex

\section{Intertemporal Consumption and Asset Pricing}

This chapter studies choice across time. The first part rewrites intertemporal consumption as an ordinary consumer-choice problem with dated consumption goods and relative prices determined by the interest rate. The second part introduces discounted utility and the Euler-style optimality condition for present versus future consumption. The final part applies present-value logic to asset pricing, net present value, perpetuities, installment loans, and asset returns under taxes.

\subsection{Intertemporal Consumption: Two-Period Model}

Consider a consumer who lives for two periods, $t=0,1$. Let
\[
x=(x_0,x_1)\in \mathbb{R}_+^2
\]
denote the consumption stream, where $x_t$ is expenditure on consumption in period $t$. Let
\[
y=(y_0,y_1)\in \mathbb{R}_+^2
\]
denote the income stream. The consumer can save or borrow between the two periods at interest rate $r\ge 0$.

\begin{definition}{Two-period savings formulation}{ch4-two-period-savings}
Let $s_0 \in \mathbb{R}$ denote period-0 saving, where
\[
s_0>0 \text{ means saving}, \qquad s_0<0 \text{ means borrowing}.
\]
The two-period budget equations are
\[
x_0+s_0=y_0,
\qquad
x_1=(1+r)s_0+y_1.
\]
\end{definition}

The first equation says current income is split between current consumption and saving. The second says future consumption equals future income plus the gross return on current saving.

\begin{theorem}{Present-value and future-value budget constraints}{ch4-pv-fv}
The two-period savings system
\[
x_0+s_0=y_0,
\qquad
x_1=(1+r)s_0+y_1
\]
is equivalent to either of the following:
\[
(1+r)x_0+x_1=(1+r)y_0+y_1,
\]
or
\[
x_0+\frac{x_1}{1+r}=y_0+\frac{y_1}{1+r}.
\]
\end{theorem}

\begin{proof}
From
\[
x_0+s_0=y_0,
\]
multiply by $1+r$:
\[
(1+r)x_0+(1+r)s_0=(1+r)y_0.
\]
Add the second budget equation:
\[
x_1=(1+r)s_0+y_1.
\]
Cancelling $(1+r)s_0$ gives
\[
(1+r)x_0+x_1=(1+r)y_0+y_1.
\]
Dividing through by $1+r$ yields
\[
x_0+\frac{x_1}{1+r}=y_0+\frac{y_1}{1+r}.
\]
\end{proof}

\begin{definition}{Present value and future value}{ch4-pv-fv-def}
For a two-period stream $(a_0,a_1)$:
\begin{itemize}
    \item its \emph{future value} at date 1 is
    \[
    FV(a_0,a_1)=(1+r)a_0+a_1;
    \]
    \item its \emph{present value} at date 0 is
    \[
    PV(a_0,a_1)=a_0+\frac{a_1}{1+r}.
    \]
\end{itemize}
Thus an intertemporal consumption plan is affordable if and only if
\[
PV(x_0,x_1)=PV(y_0,y_1).
\]
\end{definition}

This turns intertemporal choice into an ordinary consumer problem with two goods:
\[
p=(p_0,p_1)=\left(1,\frac{1}{1+r}\right),
\]
so that
\[
p\cdot x = p\cdot y.
\]
The ``price'' of future consumption is therefore
\[
p_1=\frac{1}{1+r}.
\]
A rise in $r$ lowers the price of future consumption relative to present consumption.

\begin{keyformula}[title={Key Formula: Two-Period Intertemporal Budget}]
\[
x_0+\frac{x_1}{1+r}=y_0+\frac{y_1}{1+r}.
\]
Equivalently, with
\[
p=\left(1,\frac{1}{1+r}\right),
\]
the budget set is
\[
B(p,p\cdot y)=\{x\in \mathbb{R}_+^2 : p\cdot x \le p\cdot y\}.
\]
\end{keyformula}

\begin{examplebox}[title={Worked Example: present value, future value, and saving}]
Suppose
\[
y=(190,11), \qquad r=0.1,
\]
and utility is
\[
U(x_0,x_1)=x_0x_1.
\]
The present value of income is
\[
PV(y)=190+\frac{11}{1.1}=190+10=200.
\]
The future value of income is
\[
FV(y)=1.1(190)+11=209+11=220.
\]

For Cobb--Douglas utility $U(x_0,x_1)=x_0x_1$, the optimal demand with prices
\[
\left(1,\frac{1}{1+r}\right)
\]
and present-value wealth $m$ is
\[
x(r,m)=\left(\frac{m}{2},\frac{(1+r)m}{2}\right).
\]
Hence, with $m=200$ and $r=0.1$,
\[
x(0.1,200)=\left(100,\frac{1.1\cdot 200}{2}\right)=(100,110).
\]

So the consumer saves
\[
s_0=y_0-x_0=190-100=90
\]
in period 0, and this becomes
\[
(1+r)s_0=1.1\cdot 90=99
\]
in period 1. Since future income is $11$, future consumption is
\[
99+11=110.
\]

Therefore the optimal plan is
\[
\boxed{x=(100,110), \qquad s_0=90.}
\]
\end{examplebox}

\subsection{Comparative Statics with Respect to the Interest Rate}

Because
\[
p_1=\frac{1}{1+r},
\]
a rise in the interest rate lowers the relative price of future consumption. Thus intertemporal comparative statics are a direct application of the consumer theory from Chapters 2 and 3.

\begin{definition}{Saver, borrower, and net demand for future consumption}{ch4-saver-borrower}
Given income stream $y=(y_0,y_1)$ and chosen consumption stream $x=(x_0,x_1)$:
\begin{itemize}
    \item the consumer is a \emph{saver} if
    \[
    x_1>y_1
    \quad \Longleftrightarrow \quad
    s_0>0;
    \]
    \item the consumer is a \emph{borrower} if
    \[
    x_1<y_1
    \quad \Longleftrightarrow \quad
    s_0<0.
    \]
\end{itemize}
Thus a saver is a net buyer of future consumption, while a borrower is a net seller of future consumption.
\end{definition}

This matches the endowment logic from Chapter 3 with dated consumption as the two goods and income stream $y$ as the endowment.

\begin{theorem}{Effect of an increase in the interest rate on saver/borrower status}{ch4-saver-status}
Suppose the interest rate rises from $r$ to $r'>r$.
\begin{enumerate}
    \item A saver stays a saver.
    \item A borrower may remain a borrower or may become a saver.
\end{enumerate}
\end{theorem}

\begin{proof}
Future consumption has price
\[
q(r)=\frac{1}{1+r}.
\]
If $r'>r$, then
\[
q(r')<q(r),
\]
so the price of future consumption falls. By Chapter 3, when the price of a good falls, a net buyer of that good remains a net buyer. Thus a saver stays a saver. A net seller may remain a seller or may switch to being a buyer.
\end{proof}

To apply the Slutsky decomposition, write Marshallian intertemporal demand as
\[
x(r,m)=\bigl(x_0(r,m),x_1(r,m)\bigr),
\]
where $m$ is present-value wealth. If the interest rate changes from $r$ to $r'$, then the present value of the income stream changes from
\[
m=y_0+\frac{y_1}{1+r}
\]
to
\[
m'=y_0+\frac{y_1}{1+r'}.
\]

\begin{definition}{Slutsky decomposition for intertemporal choice}{ch4-slutsky-r}
Let
\[
\Delta x_1 := x_1(r',m')-x_1(r,m)
\]
be the total change in future consumption. Define the Slutsky compensated present-value wealth at the new interest rate by
\[
m^s := x_0(r,m)+\frac{x_1(r,m)}{1+r'}.
\]
Then
\[
\Delta x_1 = \Delta x_1^s + \Delta x_1^n,
\]
where
\[
\Delta x_1^s := x_1(r',m^s)-x_1(r,m)
\]
is the substitution effect, and
\[
\Delta x_1^n := x_1(r',m')-x_1(r',m^s)
\]
is the income effect.
\end{definition}

Since $r'>r$ implies the price of future consumption falls, the substitution effect for future consumption is nonnegative:
\[
\Delta x_1^s \ge 0.
\]

If future consumption is normal and the consumer is a saver, then the income effect is also nonnegative. The reason is that a saver benefits when the interest rate rises: the price of the good she is a net buyer of falls.

\begin{theorem}{Interest-rate comparative statics for savers}{ch4-saver-theorem}
Suppose future consumption is normal and the consumer is a saver at the initial interest rate. If $r$ increases, then
\[
x_1(r',m') \ge x_1(r,m).
\]
\end{theorem}

\begin{proof}
When $r$ rises, the price of future consumption falls, so the substitution effect satisfies
\[
\Delta x_1^s \ge 0.
\]
Because the consumer is a saver, she is a net buyer of future consumption. Thus the fall in its price raises her welfare and, if future consumption is normal, the associated income effect also satisfies
\[
\Delta x_1^n \ge 0.
\]
Hence
\[
\Delta x_1=\Delta x_1^s+\Delta x_1^n\ge 0.
\]
\end{proof}

\begin{examtip}
A rise in $r$ does \emph{not} mean current wealth necessarily rises. The correct logic is:
\begin{enumerate}
    \item the price of future consumption falls;
    \item substitution pushes toward more future consumption;
    \item the welfare effect depends on whether the consumer is a saver or borrower;
    \item for a saver, both substitution and income effects push toward more future consumption if future consumption is normal.
\end{enumerate}
For a borrower, the overall effect on future consumption is ambiguous.
\end{examtip}

\subsection{Multi-Period Intertemporal Choice}

Now suppose the consumer lives for periods
\[
t=0,1,\dots,T.
\]
Let
\[
x=(x_0,x_1,\dots,x_T)\in \mathbb{R}_+^{T+1},
\qquad
y=(y_0,y_1,\dots,y_T)\in \mathbb{R}_+^{T+1}.
\]
Let $s_t$ denote assets carried from period $t$ to period $t+1$, with
\[
s_{-1}=0, \qquad s_T=0.
\]

\begin{definition}{Multi-period intertemporal budget equations}{ch4-multi-budget}
For each period $t=0,1,\dots,T$,
\[
x_t+s_t=(1+r)s_{t-1}+y_t.
\]
\end{definition}

\begin{theorem}{Multi-period present-value budget constraint}{ch4-multi-pv}
The multi-period budget system is equivalent to
\[
\sum_{t=0}^{T}\frac{x_t}{(1+r)^t}
=
\sum_{t=0}^{T}\frac{y_t}{(1+r)^t}.
\]
Equivalently, with price vector
\[
p=\left(1,\frac{1}{1+r},\frac{1}{(1+r)^2},\dots,\frac{1}{(1+r)^T}\right),
\]
the budget constraint is
\[
p\cdot x = p\cdot y.
\]
\end{theorem}

\begin{proof}
Divide each period budget equation
\[
x_t+s_t=(1+r)s_{t-1}+y_t
\]
by $(1+r)^t$:
\[
\frac{x_t}{(1+r)^t}+\frac{s_t}{(1+r)^t}
=
\frac{s_{t-1}}{(1+r)^{t-1}}+\frac{y_t}{(1+r)^t}.
\]
Summing from $t=0$ to $T$ causes the asset terms to telescope:
\[
\sum_{t=0}^{T}\frac{s_t}{(1+r)^t}
=
\sum_{t=0}^{T}\frac{s_{t-1}}{(1+r)^{t-1}}
\]
because $s_{-1}=s_T=0$. Hence
\[
\sum_{t=0}^{T}\frac{x_t}{(1+r)^t}
=
\sum_{t=0}^{T}\frac{y_t}{(1+r)^t}.
\]
\end{proof}

\begin{theorem}{Only present value matters for demand}{ch4-only-pv-matters}
Suppose the consumer solves
\[
\max_{x\in \mathbb{R}_+^{T+1}} U(x)
\quad \text{subject to} \quad
p\cdot x \le p\cdot y.
\]
If two income streams $y$ and $y'$ satisfy
\[
p\cdot y = p\cdot y',
\]
then they generate the same budget set and therefore the same demand.
\end{theorem}

\begin{proof}
The feasible set depends only on the scalar wealth level $p\cdot y$. If
\[
p\cdot y = p\cdot y',
\]
then
\[
\{x : p\cdot x \le p\cdot y\}
=
\{x : p\cdot x \le p\cdot y'\}.
\]
Hence the maximization problem is identical, so the solution is the same.
\end{proof}

\begin{examplebox}[title={Worked Example: different income timings with the same present value}]
Let $T=2$ and $r=\tfrac12$. Then
\[
p=\left(1,\frac{2}{3},\left(\frac{2}{3}\right)^2\right).
\]
Consider the three income streams
\[
y=(2,3,9), \qquad y'=(8,0,0), \qquad y''=(0,0,18).
\]
Their present values are:
\[
p\cdot y = 2+\frac{2}{3}\cdot 3 + \left(\frac{2}{3}\right)^2\cdot 9 = 2+2+4=8,
\]
\[
p\cdot y' = 8,
\]
\[
p\cdot y'' = \left(\frac{2}{3}\right)^2\cdot 18 = \frac{4}{9}\cdot 18 = 8.
\]
Thus all three income streams have the same present value. Therefore they generate the same intertemporal budget set and the same optimal demand.

So the correct conclusion is:
\[
\boxed{
(2,3,9),\ (8,0,0),\ \text{and } (0,0,18)
\text{ are equivalent for choice purposes at } r=\tfrac12.
}
\]
\end{examplebox}

\begin{commonmistake}
Do not say that the \emph{timing} of income never matters. It matters through its effect on present value. What Lecture 4 shows is that if the present value is the same, then the timing alone does not change the budget set under perfect borrowing and lending at rate $r$.
\end{commonmistake}

\subsection{Inflation and the Real Interest Rate}

Suppose the monetary price of one unit of consumption rises by factor $1+\pi$ each period. Then one unit of consumption costs:
\[
1 \text{ in period 0}, \quad 1+\pi \text{ in period 1}, \quad (1+\pi)^2 \text{ in period 2}, \dots
\]

\begin{definition}{Intertemporal budget with inflation}{ch4-inflation-budget}
With inflation rate $\pi$, the period-$t$ budget equation becomes
\[
(1+\pi)^t x_t + s_t = (1+r)s_{t-1} + (1+\pi)^t y_t.
\]
\end{definition}

Dividing by $(1+r)^t$ and summing over time yields the \emph{real} present-value budget:

\begin{theorem}{Real present-value budget}{ch4-real-pv}
With inflation rate $\pi$ and nominal interest rate $r$,
\[
\sum_{t=0}^{T}\left(\frac{1+\pi}{1+r}\right)^t x_t
=
\sum_{t=0}^{T}\left(\frac{1+\pi}{1+r}\right)^t y_t.
\]
Hence the real price vector is
\[
p=\left(1,\frac{1+\pi}{1+r},\left(\frac{1+\pi}{1+r}\right)^2,\dots,\left(\frac{1+\pi}{1+r}\right)^T\right).
\]
\end{theorem}

\begin{definition}{Real interest rate}{ch4-real-interest}
The \emph{real interest rate} $\rho$ is defined by
\[
1+\rho = \frac{1+r}{1+\pi}.
\]
Equivalently,
\[
\rho = \frac{1+r}{1+\pi}-1 = \frac{r-\pi}{1+\pi} \approx r-\pi.
\]
\end{definition}

The approximation
\[
\rho \approx r-\pi
\]
is accurate when inflation is moderate. It is the standard ``nominal minus inflation'' rule.

\begin{examplebox}[title={Worked Example: computing the real interest rate}]
Suppose
\[
r=0.32, \qquad \pi=0.20.
\]
Then
\[
1+\rho = \frac{1.32}{1.20}=1.1,
\]
so
\[
\rho = 0.1.
\]
If the income stream is
\[
y=(100,132),
\]
then its real present value is
\[
100+\frac{132}{1.1}=100+120=220.
\]
Hence
\[
\boxed{\rho=10\%, \qquad PV_{\text{real}}(y)=220.}
\]
\end{examplebox}

\subsection{Discounted Utility and Optimal Consumption Smoothing}

Lecture 4 uses the standard discounted-utility specification
\[
U(x_0,x_1)=u(x_0)+\delta u(x_1),
\]
where $u$ is the felicity function and $\delta \in [0,1]$ is the discount factor.

\begin{definition}{Discounted utility}{ch4-discounted-utility}
A two-period discounted-utility representation is
\[
U(x_0,x_1)=u(x_0)+\delta u(x_1),
\]
where:
\begin{itemize}
    \item $u:\mathbb{R}_+\to \mathbb{R}$ is increasing;
    \item $u$ is concave, so marginal utility is diminishing;
    \item $\delta \in [0,1]$ measures the value placed on future felicity.
\end{itemize}
\end{definition}

\begin{theorem}{Concavity implies diminishing marginal utility}{ch4-dmu}
If $u$ is concave and differentiable, then $u'$ is decreasing.
\end{theorem}

\begin{proof}
Concavity means that for any $x>x'$ and any $\alpha>0$,
\[
u(x+\alpha)-u(x) \le u(x'+\alpha)-u(x').
\]
Divide by $\alpha$ and let $\alpha \to 0$ to obtain
\[
u'(x) \le u'(x').
\]
Hence $u'$ is decreasing.
\end{proof}

\begin{theorem}{Euler-type first-order condition for discounted utility}{ch4-euler}
Suppose the consumer maximizes
\[
u(x_0)+\delta u(x_1)
\]
subject to
\[
x_0+\frac{x_1}{1+r}=m,
\]
and the optimum is interior. Then
\[
\frac{u'(x_0)}{\delta u'(x_1)} = 1+r,
\]
equivalently,
\[
u'(x_0) = (1+r)\delta u'(x_1).
\]
\end{theorem}

\begin{proof}
The marginal utilities are
\[
MU_0=u'(x_0), \qquad MU_1=\delta u'(x_1).
\]
By the interior tangency condition from Chapter 3,
\[
\frac{MU_0}{MU_1} = \frac{p_0}{p_1}.
\]
Here
\[
p_0=1, \qquad p_1=\frac{1}{1+r},
\]
so
\[
\frac{p_0}{p_1}=1+r.
\]
Hence
\[
\frac{u'(x_0)}{\delta u'(x_1)}=1+r,
\]
which rearranges to
\[
u'(x_0)=(1+r)\delta u'(x_1).
\]
\end{proof}

This condition compares the marginal value of consuming one more unit today with the discounted marginal value of consuming one more unit tomorrow, adjusted by the gross return from saving.

\begin{theorem}{Consumption pattern under discounted utility}{ch4-pattern}
Assume $u$ is differentiable and strictly concave, so $u'$ is strictly decreasing.
\begin{enumerate}
    \item If $(1+r)\delta = 1$, then $x_0=x_1$.
    \item If $(1+r)\delta > 1$, then $x_1>x_0$.
    \item If $(1+r)\delta < 1$, then $x_0>x_1$.
\end{enumerate}
\end{theorem}

\begin{proof}
From Theorem \ref{thm:ch4-euler},
\[
u'(x_0)=(1+r)\delta u'(x_1).
\]
If $(1+r)\delta=1$, then
\[
u'(x_0)=u'(x_1),
\]
and strict monotonicity of $u'$ implies
\[
x_0=x_1.
\]
If $(1+r)\delta>1$, then
\[
u'(x_0)>u'(x_1).
\]
Because $u'$ is decreasing, this implies
\[
x_0<x_1.
\]
Similarly, if $(1+r)\delta<1$, then
\[
u'(x_0)<u'(x_1),
\]
so
\[
x_0>x_1.
\]
\end{proof}

\begin{examplebox}[title={Worked Example: logarithmic felicity}]
Suppose
\[
u(x)=\ln x.
\]
Then
\[
u'(x)=\frac{1}{x}.
\]
The first-order condition becomes
\[
\frac{1}{x_0} = (1+r)\delta \frac{1}{x_1},
\]
so
\[
x_1=(1+r)\delta x_0.
\]

If
\[
r=0.1, \qquad \delta=0.9,
\]
then
\[
(1+r)\delta = 1.1\times 0.9 = 0.99,
\]
hence
\[
x_1=0.99x_0.
\]
So optimal future consumption is slightly below present consumption:
\[
\boxed{x_1 < x_0.}
\]
\end{examplebox}

\begin{examtip}
The single most important condition in discounted-utility questions is
\[
u'(x_0)=(1+r)\delta u'(x_1).
\]
Always interpret the product $(1+r)\delta$:
\begin{itemize}
    \item high $\delta$ means patience;
    \item high $r$ makes future consumption relatively cheap;
    \item both push consumption toward the future.
\end{itemize}
\end{examtip}

\subsection{Asset Pricing by No-Arbitrage}

An asset is a stream of future monetary payoffs. Lecture 4 prices assets by the no-arbitrage principle.

\begin{definition}{Arbitrage}{ch4-arbitrage}
An \emph{arbitrage opportunity} is a trading strategy that yields a nonnegative payoff in every contingency and a strictly positive payoff in at least one contingency, with no net cost or riskless profit at the prevailing prices.
\end{definition}

\begin{theorem}{Present-value pricing of a certain asset}{ch4-pv-asset}
Suppose an asset pays deterministic cash flow
\[
y=(y_0,y_1,\dots,y_T).
\]
If borrowing and lending are possible at interest rate $r$ and there is no arbitrage, then the asset's time-0 price must be
\[
P_0 = \sum_{t=0}^{T}\frac{y_t}{(1+r)^t}.
\]
\end{theorem}

\begin{proof}
Let
\[
PV(y)=\sum_{t=0}^{T}\frac{y_t}{(1+r)^t}.
\]

If $P_0<PV(y)$, borrow $P_0$ at time 0 and buy the asset. The asset's cash flows are sufficient to service the debt and still leave a strictly positive surplus in present-value terms, which is an arbitrage.

If $P_0>PV(y)$, short the asset, receive $P_0$ at time 0, invest the proceeds, and use the investment returns to cover the promised asset payments. The excess $P_0-PV(y)$ yields arbitrage.

Hence no-arbitrage requires
\[
P_0=PV(y).
\]
\end{proof}

\begin{definition}{Net present value}{ch4-npv}
If acquiring an asset requires payment stream
\[
m=(m_0,m_1,\dots,m_T)
\]
and yields payoff stream
\[
y=(y_0,y_1,\dots,y_T),
\]
then its \emph{net present value} is
\[
NPV=\sum_{t=0}^{T}\frac{y_t-m_t}{(1+r)^t}.
\]
\end{definition}

A project is attractive at discount rate $r$ if its net present value is positive.

\begin{theorem}{Perpetuity formula}{ch4-perpetuity}
A perpetuity that pays a fixed amount $m$ each period forever, starting in period 1, has present value
\[
PV = \frac{m}{r}.
\]
\end{theorem}

\begin{proof}
Let
\[
PV=\frac{m}{1+r}+\frac{m}{(1+r)^2}+\frac{m}{(1+r)^3}+\cdots.
\]
Multiply both sides by $\frac{1}{1+r}$:
\[
\frac{PV}{1+r}
=
\frac{m}{(1+r)^2}+\frac{m}{(1+r)^3}+\cdots.
\]
Subtract this from the original expression:
\[
PV-\frac{PV}{1+r}=\frac{m}{1+r}.
\]
So
\[
\frac{r}{1+r}PV=\frac{m}{1+r},
\]
hence
\[
PV=\frac{m}{r}.
\]
\end{proof}

\begin{examplebox}[title={Worked Example: net present value of a delayed perpetuity project}]
Suppose a project requires payment of $1$ in each of the first three periods and then pays $2$ forever from period 3 onward. Its net present value is
\[
NPV
=
-1-\frac{1}{1+r}-\frac{1}{(1+r)^2}
+\frac{2}{(1+r)^3}+\frac{2}{(1+r)^4}+\cdots.
\]
The infinite tail is a perpetuity starting at period 3, so
\[
\frac{2}{(1+r)^3}+\frac{2}{(1+r)^4}+\cdots
=
\frac{1}{(1+r)^2}\left(\frac{2}{1+r}+\frac{2}{(1+r)^2}+\cdots\right)
=
\frac{1}{(1+r)^2}\cdot \frac{2}{r}.
\]
Therefore
\[
NPV
=
\frac{2}{r(1+r)^2}
-
\left(
1+\frac{1}{1+r}+\frac{1}{(1+r)^2}
\right).
\]
So the correct compact expression is
\[
\boxed{
NPV
=
\frac{2}{r(1+r)^2}
-
\left(
1+\frac{1}{1+r}+\frac{1}{(1+r)^2}
\right).
}
\]
\end{examplebox}

\begin{examplebox}[title={Worked Example: installment loan interest rate}]
An installment loan gives $m=1000$ upfront and requires monthly payments of $y=100$ for $T=12$ months. The correct monthly interest rate $r$ solves
\[
1000
=
\frac{100}{1+r}
+
\frac{100}{(1+r)^2}
+\cdots+
\frac{100}{(1+r)^{12}}.
\]
This is \emph{not} the same as taking
\[
\frac{12\cdot 100}{1000}-1=20\%.
\]
That simple ratio ignores the fact that payments are made over time rather than all at the end.

Numerically, the solution is approximately
\[
r \approx 0.03
\]
per month, i.e. about $3\%$ monthly. The implied annual rate is therefore much larger than $20\%$; Lecture 4 reports it as roughly $41\%$ per year.

The exam logic is:
\begin{enumerate}
    \item write the present-value equation for the loan;
    \item solve for the periodic interest rate;
    \item do not treat the total repayment as if it occurred only in the final period.
\end{enumerate}
\end{examplebox}

\subsection{Asset Price Changes and Rates of Return}

Suppose an asset does not pay dividends in period 0 and its price rises from $p_0$ to $p_1$ with certainty by next period.

\begin{theorem}{No-arbitrage return condition}{ch4-return}
If the asset pays no interim cash flow and can be traded freely, then no arbitrage implies
\[
r=\frac{p_1-p_0}{p_0}.
\]
Equivalently,
\[
p_0=\frac{p_1}{1+r}.
\]
\end{theorem}

\begin{proof}
If
\[
\frac{p_1-p_0}{p_0}>r,
\]
borrow $p_0$, buy the asset, sell it next period for $p_1$, and repay $(1+r)p_0$. The surplus
\[
p_1-(1+r)p_0>0
\]
is an arbitrage.

If
\[
\frac{p_1-p_0}{p_0}<r,
\]
short the asset, receive $p_0$, invest at rate $r$, and next period buy back the asset for $p_1$. The surplus
\[
(1+r)p_0-p_1>0
\]
is an arbitrage.

Therefore only
\[
\frac{p_1-p_0}{p_0}=r
\]
can hold.
\end{proof}

\begin{policynote}[title={Policy Application: housing as an asset}]
A house pays housing services rather than direct cash. If the monetary rental value of current housing services is $y_0$ and next period's house price is $p_1$, then no-arbitrage implies
\[
p_0 = y_0 + \frac{p_1}{1+r}.
\]
So the value of a house is the value of current housing services plus the discounted resale value.
\end{policynote}

\begin{policynote}[title={Policy Application: depletable resources}]
If a depletable resource such as oil yields no extraction cost and must eventually be replaced by a perfect substitute with cost $C$ at terminal date $T$, then no-arbitrage implies
\[
p_t = \frac{p_{t+1}}{1+r}.
\]
Hence
\[
p_0 = \frac{p_T}{(1+r)^T} = \frac{C}{(1+r)^T}.
\]
This is the basic present-value logic behind the price path of a scarce exhaustible resource.
\end{policynote}

\subsection{Taxes and Asset Returns}

Lecture 4 also studies taxes on interest income and asset returns.

\begin{definition}{After-tax interest rate}{ch4-after-tax}
If the tax rate on interest income is $\tau$, then one dollar saved becomes
\[
1+r(1-\tau)
\]
next period. Thus the effective after-tax interest rate is
\[
r(1-\tau).
\]
\end{definition}

This means taxes reduce the reward to saving and soften the incentive to transfer consumption into the future.

Now consider two assets:
\begin{itemize}
    \item asset $A$ is untaxed;
    \item asset $B$ faces tax rate $\tau$ on its capital gain.
\end{itemize}

\begin{theorem}{Required pre-tax return on a taxed asset}{ch4-taxed-asset}
Suppose asset $B$ has price $p_0^B$ today and price $p_1^B$ next period, and the tax rate on its return is $\tau$. No arbitrage implies
\[
p_0^B = \left(\frac{1-\tau}{1-\tau+r}\right)p_1^B,
\]
equivalently,
\[
\frac{p_1^B-p_0^B}{p_0^B}=\frac{r}{1-\tau}.
\]
\end{theorem}

\begin{proof}
If an investor buys asset $B$ at time 0 for $p_0^B$, then next period the gross proceeds net of tax are
\[
p_1^B - \tau(p_1^B-p_0^B).
\]
No arbitrage requires this to equal the gross return from investing $p_0^B$ at interest rate $r$:
\[
(1+r)p_0^B = p_1^B - \tau(p_1^B-p_0^B).
\]
Rearranging:
\[
(1+r)p_0^B = (1-\tau)p_1^B + \tau p_0^B,
\]
so
\[
(1-\tau+r)p_0^B = (1-\tau)p_1^B.
\]
Hence
\[
p_0^B = \left(\frac{1-\tau}{1-\tau+r}\right)p_1^B.
\]
Subtract $p_0^B$ from both sides of the equivalent expression
\[
(1-\tau+r)p_0^B=(1-\tau)p_1^B
\]
to get
\[
rp_0^B = (1-\tau)(p_1^B-p_0^B),
\]
which gives
\[
\frac{p_1^B-p_0^B}{p_0^B}=\frac{r}{1-\tau}.
\]
\end{proof}

Thus the taxed asset must offer a higher \emph{pre-tax} return than the untaxed asset in order to remain competitive after tax.

\subsection{Chapter Summary}

\begin{keyformula}[title={Key Formula: Lecture 4 Summary}]
\begin{align*}
x_0+s_0 &= y_0, \\
x_1 &= (1+r)s_0+y_1, \\
x_0+\frac{x_1}{1+r} &= y_0+\frac{y_1}{1+r}, \\
p &= \left(1,\frac{1}{1+r}\right), \\
\sum_{t=0}^{T}\frac{x_t}{(1+r)^t} &= \sum_{t=0}^{T}\frac{y_t}{(1+r)^t}, \\
1+\rho &= \frac{1+r}{1+\pi}, \\
U(x_0,x_1) &= u(x_0)+\delta u(x_1), \\
u'(x_0) &= (1+r)\delta u'(x_1), \\
P_0 &= \sum_{t=0}^{T}\frac{y_t}{(1+r)^t}, \\
NPV &= \sum_{t=0}^{T}\frac{y_t-m_t}{(1+r)^t}, \\
PV_{\text{perpetuity}} &= \frac{m}{r}, \\
r &= \frac{p_1-p_0}{p_0}, \\
p_0^B &= \left(\frac{1-\tau}{1-\tau+r}\right)p_1^B.
\end{align*}
\end{keyformula}

\begin{examtip}
The highest-yield skills from Lecture 4 are:
\begin{itemize}
    \item convert intertemporal budgets into present-value form;
    \item interpret $1/(1+r)$ as the price of future consumption;
    \item distinguish savers from borrowers and use price-change logic correctly;
    \item use
    \[
    u'(x_0)=(1+r)\delta u'(x_1)
    \]
    for discounted-utility problems;
    \item compute present values, real interest rates, and net present values cleanly;
    \item apply no-arbitrage to price assets and derive required returns;
    \item adjust return formulas correctly when taxes are imposed.
\end{itemize}
\end{examtip}